\section{Principle X: Tradition}

\begin{quotex}
Maurras concludes his list of fundamental principles with Tradition. This will be a disappointment to those who want to reduce everything to the rational, since traditional is beyond that. It is tied, as he says, to ``blood and soil", that is, one's ancestors and the heritage they bequeathed. The deracinated modern mind is limited by the idea, the simulacra of the ``proposition" nation.

Yet, tradition is not blind. It successes serve as the paradigm for the future, it failures as lessons to be learned; there is no point to perpetuate errors. This is a living tradition, not stuck in the past through inertia. A true Tradition is never lost, it is only hidden. If it had one starting point, it can have a second. It will being anew only by men who know who they are. 

\end{quotex}
Tradition means handing down.

Tradition gathers the forces of blood and soil. It is retained even when leaving one's country, as an eternal temptation to return to it.

True tradition is critical, and for lack of its distinctions, the past is no longer of use for anything, its successes stop being models, its setbacks stop being lessons.

In all tradition, as in all heritage, a reasonable being deducts and must deduct the liabilities.

Tradition is not inertia, its opposite; heritage is not nepotism, its counterfeit.

All traditions had a beginning and the sentiments of monarchistic fidelity, if they rise up very high, do not do so indefinitely: what began can begin again; what had one point of departure can find a second.

The opposition between reason and tradition, which is the same as the antithesis of reality and the idea, or art and nature, and that can be assimilated to the opposition of vinegar and oil, the sweet and the bitter, the fluid and the solid, in a cosmogony of infant peoples.



\flrightit{Posted on 2012-12-23 by Cologero }

\begin{center}* * *\end{center}

\begin{footnotesize}\begin{sffamily}



\texttt{White Stag Knight on 2012-12-23 at 11:31 said: }

Maurras, must be handled delicately and with dialectic quasi-Rabbinic intellectualism of finesse. He was not condemned by Catholicism without reason. He cannot be the center of our Right Hand Path counter-revolution. He can be but a source of truths to be mined, but no Doctor. 

The nationalism of Maurras, in fact, signified in one sense the miasma of international European aristocracy and the fragmentation of Christendom.

We aim to resurrect the Holy Roman Imperium of Christ Rex. 

No single nation can exist in self-deifying separatism and claim special privileges.

The Empire is Universal.


\hfill

\texttt{Cologero on 2012-12-23 at 12:32 said: }

Obviously, White Knight, we can't disagree with you. However, using his own standard, if we add Maurras' assets and deduct his liabilities, we are still in the black. My prime motivation has been this: to show that a pagan non-theist relying on positive science can still discern Tradition. That is because the Logos is made visible in creation. There are many men today in that position who would still nevertheless want to align with Tradition. Maurras shows that honest science does not necessarily have that corrosive effect assumed by the anti-Tradition.

His hyper-nationalism is not the ultimate reality, but remember that the Empire does not deny nor annihilate nationhood. At one point, there was a group around Maurras that floated the idea of an alliance of the Latin-speaking nations. That would have been a start toward reunifying the ecumene. Had that materialized, the role of Romania would have been important in uniting the East and the West.


\hfill

\texttt{Dominion on 2012-12-23 at 14:46 said: }

Thank you so much for translating this work. I'm beginning to think it would be worth it to learn French simply to read the originals. Seeing as the Action Francaise was of a royalist nature, perhaps there is room here for an article comparing Maurras' view of the royal vs the religious spheres in the social order with those of Evola and Guenon? I would guess that he would approximate Evola to a greater degree, having the Church being more subject to the King, but if there is one thing I have learned from these principles it is that he is not to be underestimated in his depth. 

With regards to what you say about the role of nationhood in Tradition, I see the proper view more and more espoused in the ``Identitaire" movement that is currently exploding over Europe. They hold the local (city and region), national, and European levels of identity as valid and complementary. I have found (observing it from a distance from across the pond) that there are many in that movement who are familiar with some Traditional works, Evola figuring large (with his shortcomings compared to Guenon or some others causing some problems, admittedly). Do you think a possibility exists of leading those who are open from a ``mere" nationalism to one informed by Tradition?


\hfill

\texttt{Jason-Adam on 2012-12-27 at 00:38 said: }

Cologero – can you provide some links in English or French explaining in more detail Maurras' plan of reuniting the Latin nations ? I was unaware that he went beyond French nationalism and would like to read more on this.


\hfill

\texttt{Cologero on 2012-12-27 at 06:08 said: }

Here is an overview: L'union des latins\footnote{\url{http://maurras.net/2011/10/06/union-des-latins/}}.

In Maurras' own words: Les Forces latines\footnote{\url{http://maurras.net/textes/196.html}}.


\hfill

\texttt{Jason-Adam on 2012-12-27 at 17:41 said: }

Thank you, Sir. I've just read both and find myself in agreement with Maurras – one point he made that I never considered before is that the southern Catholic Germans are actually descended from Roman colonists, and the failed Holy Roman Empire was an attempt to Latinise the Germans. I'm finding myself in accord with his notion of Catholicism being part of the Latin race's blood, his version of the continuation of Romanity rings truer than Evola's idea of Germany being the successor to Roma.


\end{sffamily}\end{footnotesize}
